% Acronyms
\newacronym[type=main,description={Acronimo di Asynchronous Server Gateway Interface: specifica di interfaccia per applicazioni web asincrone che definisce un contratto standard tra server e applicazioni, abilitando runtime basati su event-loop, connessioni persistenti e gestione efficiente di I/O concorrente; separa lo stack di trasporto dal codice applicativo e consente server/worker intercambiabili.}]{asgi}{ASGI}{Asynchronous Server Gateway Interface}
\newacronym[type=main,description={Acronimo di Application Program Interface: interfaccia applicativa che espone funzionalit\`a e dati di un sistema software tramite contratti stabili; consente a componenti esterni di invocare operazioni e scambiare payload strutturati senza conoscere l'implementazione interna. Include aspetti di versioning, sicurezza e documentazione per garantire integrazioni affidabili.}]{api}{API}{Application Program Interface}
\newacronym{csv}{CSV}{Comma-Separated Values}
\newacronym{html}{HTML}{HyperText Markup Language}
\newacronym{http}{HTTP}{Hypertext Transfer Protocol}
\newacronym{ide}{IDE}{Integrated Development Environment}
\newacronym{json}{JSON}{JavaScript Object Notation}
\newacronym{jwks}{JWKS}{JSON Web Key Set}
\newacronym{jwt}{JWT}{JSON Web Token}
\newacronym{gpu}{GPU}{Graphics Processing Unit}
\newacronym{crud}{CRUD}{Create, Read, Update, Delete}
\newacronym[type=main,description={Acronimo di Large Language Model: modello di intelligenza artificiale di grandi dimensioni addestrato su corpora estesi per generare, riassumere e comprendere linguaggio naturale; utilizza architetture neurali con molti parametri per produrre risposte contestuali e coerenti in numerosi domini applicativi.}]{llm}{LLM}{Large Language Model}
\newacronym[type=main,description={Acronimo di Enterprise Resource Planning: sistema informativo integrato che centralizza anagrafiche, processi amministrativi, logistici e produttivi in un unico database e set di moduli; garantisce coerenza dei dati, tracciabilit\`a delle operazioni e automazione di workflow aziendali.}]{erp}{ERP}{Enterprise Resource Planning}
\newacronym{mape}{MAPE}{Mean Absolute Percentage Error}
\newacronym{mrp}{MRP}{Material Requirements Planning}
\newacronym{rest}{REST}{Representational State Transfer}
\newacronym{rmse}{RMSE}{Root Mean Squared Error}
\newacronym{shap}{SHAP}{SHapley Additive exPlanations}
\newacronym{smape}{SMAPE}{Symmetric Mean Absolute Percentage Error}
\newacronym{uuid}{UUID}{Universally Unique Identifier}
\newacronym[type=main,description={Acronimo di Explainable Artificial Intelligence: disciplina che rende interpretabili i modelli di AI, fornendo motivazioni leggibili a supporto di decisioni automatiche; comprende tecniche locali e globali per evidenziare contributi delle variabili, verificare bias, soddisfare requisiti normativi e aumentare la fiducia degli utenti.}]{xai}{XAI}{Explainable Artificial Intelligence}
\newacronym{sdk}{SDK}{Software Development Kit}
\newacronym{uml}{UML}{Unified Modeling Language}
\newacronym{wape}{WAPE}{Weighted Absolute Percentage Error}
\newacronym{tsa}{TSA}{Termine solo acronimo}
\newacronym{xml}{XML}{eXtensible Markup Language}
\newglossaryentry{ml}{
    name={Machine Learning},
    sort=machine learning,
    description={Branca dell’intelligenza artificiale che sviluppa modelli in grado di apprendere da dati e migliorare le proprie prestazioni senza programmazione esplicita, tramite addestramento su esempi e ottimizzazione di una funzione obiettivo}
}
\newglossaryentry{continuouseintegration}{
    name={Integrazione continua},
    sort=integrazione continua,
    description={Pratica di sviluppo che prevede l’integrazione frequente del codice in un \textit{repository} condiviso, con build e test automatici a ogni modifica per rilevare errori rapidamente}
}
\newglossaryentry{productowner}{
    name={Product Owner},
    sort=product owner,
    description={Figura del framework Scrum responsabile della massimizzazione del valore del prodotto: definisce e priorizza il backlog, chiarisce obiettivi e criteri di accettazione, rappresentando le esigenze degli stakeholder}
}
\newglossaryentry{scrum}{
    name={Scrum},
    sort=scrum,
    description={Framework agile per la gestione di prodotti complessi, basato su iterazioni time-boxed (Sprint), ruoli definiti (Product Owner, Scrum Master, Developers) ed eventi cadenzati (Sprint Planning, Daily Scrum, Sprint Review, Sprint Retrospective) descritti dalla Scrum Guide (\url{https://scrumguides.org/scrum-guide.html})}
}
\newglossaryentry{scrummaster}{
    name={Scrum Master},
    sort=scrum master,
    description={Facilitatore del framework Scrum: guida il team nell’applicazione delle pratiche Scrum, rimuove impedimenti, facilita le cerimonie e promuove il miglioramento continuo}
}
\newglossaryentry{mergerequest}{
    name={Merge \textit{request}},
    sort=merge request,
    description={Richiesta di integrazione di un branch in un altro (tipicamente verso main/develop) su piattaforme di version control: include diff, discussioni, revisioni e controlli automatizzati prima del merge}
}
\newglossaryentry{overfitting}{
    name={Overfitting},
    sort=overfitting,
    description={Situazione in cui un modello apprende troppo dai dati di training, catturando rumore o peculiarità non generalizzabili e performando peggio su dati nuovi}
}
\newglossaryentry{underfitting}{
    name={Underfitting},
    sort=underfitting,
    description={Situazione in cui un modello è troppo semplice o poco allenato per catturare la struttura dei dati, mostrando bassa accuratezza sia su training sia su test}
}
\newglossaryentry{ioc}{
    name={Inversion of Control},
    sort=inversion of control,
    description={Principio architetturale in cui il flusso di controllo è invertito rispetto alla programmazione tradizionale: i componenti dichiarano dipendenze che vengono fornite da un contenitore esterno, anziché crearle direttamente, favorendo disaccoppiamento e testabilità}
}
\newglossaryentry{di}{
    name={Dependency Injection},
    sort=dependency injection,
    description={Pattern che implementa l’Inversion of Control passando dall’esterno le dipendenze necessarie a un componente (es. tramite costruttore o setter), evitando la creazione interna e permettendo configurazione e sostituzione}
}
\newglossaryentry{branch}{
    name={Branch},
    sort=branch,
    description={Ramo di sviluppo separato in un sistema di version control distribuito, usato per lavorare su modifiche isolate prima di integrararle nel ramo principale}
}

% Glossary
\newglossaryentry{erpg}{
    name={ERP},
    text={Enterprise Resource Planning},
    sort=erp,
    description={Un sistema Enterprise Resource Planning (ERP) integra in un’unica piattaforma i principali processi aziendali, come amministrazione, vendite, acquisti e magazzino, fornendo dati condivisi e coordinamento tra i reparti}
}

\newglossaryentry{sdkg}{
    name={SDK},
    text={Software Development Kit},
    sort=sdk,
    description={A Software Development Kit (SDK) is a collection of development tools in one installable package, facilitating application creation by providing a compiler, debugger, and sometimes a software framework. SDKs are typically specific to a hardware platform and operating system combination. Many application developers use specific SDKs to enable advanced functionalities such as advertisements, push notifications, etc}
}

\newglossaryentry{umlg}{
    name={UML},
    text={Unified Modeling Language},
    sort=uml,
    description={In software engineering, Unified Modeling Language (UML) is a modeling and specification language based on the object-oriented paradigm. UML serves as a "lingua franca" in the object-oriented design and programming community. Much of the industry literature uses UML to describe analytical and design solutions in a concise and understandable way for a broad audience}
}

\newglossaryentry{TermineSenzaAcronimo}{
    name={Nome del termine},
    sort=termine senza acronimo,
    description={Descrizione}
}

\newglossaryentry{jobasync}{
    name={Job asincrono},
    sort=job asincrono,
    description={Unit\`a di lavoro eseguita in background con stato tracciato (in coda, in esecuzione, completata, fallita) e consultabile tramite identificativo}
}

\newglossaryentry{callback}{
    name={Callback / Webhook},
    sort=callback,
    description={Chiamata HTTP inviata da un sistema a un endpoint esterno per notificare eventi o risultati, spesso firmata/autenticata e usata per integrazioni push}
}
\newglossaryentry{similaritysearch}{
    name={Similarity search},
    text={similarity search},
    sort=similarity search,
    description={Tecnica di ricerca che, dato un oggetto di query, restituisce gli oggetti più simili in base a una misura di distanza (es. cosine, L2) calcolata su rappresentazioni vettoriali come gli embedding. Usata per ricerca semantica, nearest neighbors e raccomandazione}
}
\newglossaryentry{dto}{
    name={DTO},
    text={DTO},
    sort=dto,
    description={Data Transfer Object: oggetto leggero e tipizzato usato per trasportare dati tra livelli o servizi, privo di logica di dominio; riduce l’accoppiamento e rende espliciti contratti di input/output}
}
\newglossaryentry{devops}{
    name={DevOps},
    text={DevOps},
    sort=devops,
    description={Insieme di pratiche e cultura collaborativa che integra sviluppo (Dev) e operations (Ops) per automatizzare build, test, rilascio, monitoraggio e feedback continuo; punta a cicli di consegna rapidi e affidabili}
}
\newglossaryentry{olap}{
    name={OLAP},
    text={OLAP},
    sort=olap,
    description={Online Analytical Processing: modello e insieme di tecniche per analisi di dati multidimensionali, tipicamente ottimizzate per interrogazioni aggregate (slice/dice) su volumi elevati; privilegia throughput e scansioni a colonne rispetto alle operazioni transazionali OLTP}
}

\newglossaryentry{featureengineering}{
    name={Feature engineering},
    sort=feature engineering,
    description={Processo di creazione, trasformazione e selezione di variabili derivate dai dati grezzi per migliorare l’addestramento e le prestazioni dei modelli}
}

\newglossaryentry{holdout}{
    name={Hold-out},
    sort=holdout,
    description={Sottinsieme dei dati tenuto da parte rispetto al training per valutare in modo indipendente le prestazioni di un modello}
}

\newglossaryentry{embedding}{
    name={Embedding},
    sort=embedding,
    description={Rappresentazione vettoriale densa di testo o altri oggetti che preserva relazioni semantiche, usata per similarity search e modelli downstream}
}

\newglossaryentry{repositorypattern}{
    name={Repository Pattern},
    sort=repository pattern,
    description={Pattern architetturale che incapsula l’accesso ai dati dietro interfacce dedicate, separando la logica di dominio dai dettagli di persistenza}
}

\newglossaryentry{usecase}{
    name={Use case},
    sort=use case,
    description={Descrizione strutturata di un’interazione utente-sistema con attori, precondizioni, flusso principale e varianti; utilizzata per raccogliere e verificare requisiti}
}

\newglossaryentry{jobrun}{
    name={Job / Run di previsione},
    sort=job run,
    description={Esecuzione identificata (spesso con UUID) che traccia uno specifico ciclo di elaborazione o previsione, con risultati, metriche e log associati}
}

\newglossaryentry{forecastpipeline}{
    name={Pipeline di previsione},
    sort=pipeline previsione,
    description={Sequenza orchestrata di fasi (acquisizione dati, preparazione feature, inferenza, spiegazioni, persistenza, notifiche) per produrre e distribuire previsioni}
}

\newglossaryentry{happypath}{
    name={Happy path},
    sort=happy path,
    description={Percorso di esecuzione ideale senza errori o eccezioni, usato nei test per validare il comportamento atteso in condizioni nominali}
}

\newglossaryentry{branchingmodel}{
    name={Modello di branching},
    sort=branching model,
    description={Strategia di gestione dei rami Git (es. main/develop/feature/fix/docs) che definisce come organizzare sviluppo, integrazione e rilascio}
}

\newglossaryentry{kanban}{
    name={Kanban board},
    sort=kanban,
    description={Lavagna organizzata in colonne (es. Backlog, Todo, In Progress, Done) per visualizzare e limitare il lavoro in corso secondo i principi Kanban}
}

\newglossaryentry{sprint}{
    name={Sprint},
    sort=sprint,
    description={Iterazione time-boxed tipica di Scrum in cui un team pianifica, realizza e rilascia un incremento di prodotto potenzialmente utilizzabile}
}

\newglossaryentry{backlog}{
    name={Backlog},
    sort=backlog,
    description={Elenco prioritizzato di attivit\`a, requisiti o user story da eseguire, gestito e raffinato nel tempo dal team di prodotto}
}
